% Figure 4: Corralling Semantic Analysis
% This file contains the figure caption and mathematical framework for Figure 4

\begin{figure*}[t]
\centering
\includegraphics[width=0.95\textwidth]{experiments_v1/04_figure/results/figure3_corralling_semantic_analysis.png}
\caption{%
\textbf{Long-Term Adaptation: Corralling Unlearns Suboptimal Semantic Transfer.}
%
\textbf{(Left)} Semantic structure of LMSYS Chat-1M dataset (N=594k prompts) projected onto first two principal components. 
The Easy cluster (PC1 $< 0.3$, blue, 94.2\% of prompts) exhibits high semantic density. 
The Hard cluster (PC1 $\geq 0.3$, red, 5.8\% of prompts) is more sparse.
%
\textbf{(Right)} Expert weight evolution during Corralling training on N=1,121 labeled samples (3 models spanning cost tiers: Mixtral \$0.27/1M, GPT-4-Turbo \$10/1M, GPT-4o \$2.50/1M) with \textbf{aggressive} learning rate ($\eta=5.0$, $\gamma=0.10$).
GPT-4o is initialized via \textbf{semantic transfer} from GPT-4-Turbo ($\gamma=0.05$), providing an informative prior.
The algorithm starts with uniform expert weights (0.5, 0.5) and rapidly shifts weight from warmup expert (orange) to tabula rasa expert (green).
Final weights: Warmup=$1.41 \times 10^{-128}$ (essentially 0), Tabula Rasa=1.000, demonstrating \textbf{complete unlearning} of warmup priors including semantic transfer.
This validates Corralling's long-term safety: semantic transfer provides short-term benefit (breaking symmetry), but with sufficient data and aggressive learning, the system adapts beyond initial priors to discover the optimal policy.
Performance: Cumulative regret = 59.33 $\pm$ 3.40, Average reward = 0.939 $\pm$ 0.003 (mean $\pm$ std across 3 seeds).
Learned policy: GPT-4o (70.8\%), Mixtral (23.2\%), GPT-4-Turbo (6.0\%), achieving \$2.43/1M tokens (75.7\% cost reduction).
\textbf{Connection to Figure~\ref{fig:accelerated-adoption}:} This experiment uses 50$\times$ higher learning rate than Figure~\ref{fig:accelerated-adoption}, demonstrating that semantic transfer benefits are transient with aggressive learning. The unlearning validates system robustness: the router is not locked into potentially incorrect semantic priors.
}
\label{fig:corralling_semantic}
\end{figure*}


% ============================================================================
% Mathematical Framework: Corralling Algorithm
% ============================================================================

\subsection{Corralling: Adaptive Expert Combination}
\label{sec:corralling}

We employ a simplified variant of the \emph{Corralling} algorithm~\cite{agarwal2017corralling} to adaptively combine two bandit strategies:
a warmup expert initialized with priors from RouteLLM, and a tabula rasa expert learning from scratch.
This provides adaptive robustness: if warmup priors are harmful due to domain mismatch, the algorithm automatically shifts weight to the tabula rasa expert.

\paragraph{Problem Setup.}
We have $E$ experts (in our case, $E=2$), each implementing a bandit algorithm.
At each round $t$:
\begin{enumerate}[nosep]
    \item The meta-algorithm samples an expert $e_t \sim p_t$, where $p_t$ is a probability distribution over experts
    \item Expert $e_t$ selects a model $a_t$ based on context $x_t$
    \item We observe reward $r_t \in [0,1]$ for the selected model
    \item The meta-algorithm updates expert weights based on observed losses
\end{enumerate}

\paragraph{Importance-Weighted Loss Estimation.}
The key challenge is that we only observe the outcome for the \emph{selected} expert, not for all experts (no counterfactuals).
To obtain unbiased loss estimates, we use importance weighting:
\begin{equation}
\hat{\ell}_{t,e} = \begin{cases}
\frac{1 - r_t}{\rho_{t,e}} & \text{if } e = e_t \text{ (chosen expert)} \\
0 & \text{otherwise}
\end{cases}
\label{eq:importance_weighted_loss}
\end{equation}
where $\rho_{t,e} = p_t(e)$ is the probability of selecting expert $e$ at round $t$.

This estimator is unbiased:
\begin{equation}
\mathbb{E}_{e_t \sim p_t}[\hat{\ell}_{t,e}] 
= \sum_{e'=1}^{E} p_t(e') \cdot \frac{\mathbb{1}_{e'=e} \cdot (1-r_t)}{p_t(e')} 
= 1 - r_t
\end{equation}

\paragraph{Exponential Weights Update.}
We maintain cumulative losses $L_{t,e} = \sum_{s=1}^{t} \hat{\ell}_{s,e}$ for each expert and update weights using:
\begin{equation}
w_{t+1,e} = \frac{\exp(-\eta \cdot L_{t,e})}{\sum_{e'=1}^{E} \exp(-\eta \cdot L_{t,e'})}
\label{eq:exponential_weights}
\end{equation}
where $\eta > 0$ is the learning rate controlling adaptation speed.

\paragraph{Theoretical Motivation.}
Our architecture is grounded in the \emph{Corralling} framework~\cite{agarwal2017corralling}, which theoretically establishes that a meta-learner can achieve sublinear regret relative to the best expert in hindsight:
\begin{equation}
\mathbb{E}[\text{Regret}] \leq \min_{e \in [E]} \mathbb{E}[\text{Regret}_e] + O\left(\sqrt{T \log E}\right)
\end{equation}
While our implementation simplifies the update rule (updating only the selected expert to ensure $O(1)$ latency rather than $O(K)$), we retain the core \textbf{Importance-Weighted Loss} objective (Eq.~\ref{eq:importance_weighted_loss}).
This ensures that the system inherits the \emph{behavioral properties} of the theoretical algorithm:
\begin{itemize}[nosep]
    \item If warmup priors are helpful, the meta-learner aligns with the Warmup expert.
    \item If warmup priors are harmful (negative transfer), the importance-weighted penalty explodes, forcing a shift to the Tabula Rasa expert.
\end{itemize}

\paragraph{Technical Fairness and Expert Starvation.}
A known limitation of the Exp3-style estimator (Eq.~\ref{eq:importance_weighted_loss}) in our partial-feedback setting is the potential for ``Expert Starvation.'' Since unselected experts receive $\ell=0$ and no internal state update ($\Delta A = 0$), a Tabula Rasa expert that fails to gain traction early may effectively freeze in a high-entropy state.

We mitigate this risk through two mechanisms:
\begin{enumerate}[nosep]
    \item \textbf{Uniform Initialization:} We strictly initialize trust weights at parity ($\pi_0 = [0.5, 0.5]$). As shown in Figure~\ref{fig:corralling_semantic}, this ensures the Tabula Rasa expert receives $\sim$50\% of traffic during the critical burn-in phase ($t < 50$), allowing it to build a sufficient initial covariance model before trust differentiation begins.
    \item \textbf{Incumbent Failure Logic:} Starvation is only a risk if the dominant expert (Warmup) continues to succeed. If the Warmup expert performs poorly (as in our Domain Mismatch scenario), its cumulative loss $L_{\text{warmup}}$ rises. Because probabilities are softmax-normalized (Eq.~\ref{eq:exponential_weights}), a rise in the incumbent's loss mathematically forces an increase in the starved expert's probability, naturally breaking the starvation loop.
\end{enumerate}

Thus, our architecture favors the incumbent only as long as it remains optimal; any degradation in the Warmup prior automatically flushes traffic to the Tabula Rasa candidate.

\paragraph{Implementation Details.}
We set $\eta = 5.0$ (learning rate) and $\gamma = 0.10$ (mixing parameter), selected via systematic ablation studies over 45 configurations (5 learning rates $\times$ 3 mixing parameters $\times$ 3 random seeds).
All experts are initialized with uniform weights $w_{1,e} = 1/E$ to ensure fair initial exploration.
The warmup expert uses gamma-scaled priors ($\gamma=0.05$) from RouteLLM, while the tabula rasa expert starts from scratch with $A = I$ and $b = 0$.
For models not present in the warmup priors (GPT-4o), we employ semantic transfer: initializing from the closest available model (GPT-4-Turbo) with gamma scaling to provide a weak prior while allowing rapid adaptation through online learning.

\paragraph{Experimental Results: Long-Term Adaptation Beyond Semantic Transfer.}
Figure~\ref{fig:corralling_semantic} demonstrates Corralling's long-term adaptive behavior with \textbf{complete unlearning} (100\% weight transfer) to the tabula rasa expert.
Critically, the warmup expert includes \textbf{semantic transfer} for GPT-4o (inherited from GPT-4-Turbo with $\gamma=0.05$), yet this prior is ultimately rejected.

Analysis reveals three key findings:

\textbf{(1) Semantic Transfer Provides Short-Term Benefit:} During the first $\sim$100--200 steps, semantic transfer breaks symmetry and accelerates exploration. GPT-4o inherits GPT-4-Turbo's learned preferences, providing an informative starting point.

\textbf{(2) Long-Term Adaptation Unlearns Priors:} With aggressive learning rate ($\eta=5.0$), the meta-learner detects that warmup priors (including semantic transfer) are suboptimal. By step $\sim$300, expert weights have shifted dramatically from [0.5, 0.5] to [$\sim$0.02, $\sim$0.98]. Complete unlearning occurs by step 1,121.

\textbf{(3) Optimal Policy Differs from Priors:} The tabula rasa expert learns that:
\begin{itemize}[nosep]
    \item GPT-4o achieves comparable quality to GPT-4-Turbo at \$2.50 vs \$10.00 per 1M tokens (4$\times$ cost reduction)
    \item The \emph{Easy cluster} (94.2\% of prompts) can be served effectively by mid-tier models
    \item Mixtral handles 23.2\% of routine tasks at minimal cost (\$0.27/1M)
\end{itemize}

Notably, the algorithm optimizes purely for quality ($\lambda_{\text{cost}} = 0$); the 75.7\% cost reduction emerges naturally as the learned policy discovers more efficient model choices.

Final expert weights after training on N=1,121 samples:
\begin{itemize}[nosep]
    \item Warmup Expert (with semantic transfer): $1.41 \times 10^{-128}$ (essentially 0\%)
    \item Tabula Rasa Expert: 1.000 (100.0\%)
\end{itemize}

\textbf{Interpretation:} The complete unlearning validates two critical properties:
\begin{enumerate}[nosep]
    \item \textbf{System Robustness:} The router is not locked into potentially incorrect semantic priors. Even when semantic transfer provides an informative initialization, the system adapts to the true data distribution.
    \item \textbf{Mechanism Clarification:} Semantic transfer benefits come from \emph{implicit regularization} (breaking symmetry, providing initial variance) rather than accurate performance prediction. With sufficient data, the system converges to the optimal policy regardless of prior quality.
\end{enumerate}

The system achieved cumulative regret of 59.33 $\pm$ 3.40 and average reward of 0.939 $\pm$ 0.003 (mean $\pm$ std across 3 random seeds), confirming high performance with consistent reproducibility despite unlearning the initial priors.

\paragraph{Hyperparameter Optimization.}
To ensure robust performance, we conducted systematic ablation studies across 45 configurations (5 learning rates $\in \{0.1, 0.5, 1.0, 2.0, 5.0\}$ $\times$ 3 mixing parameters $\in \{0.0, 0.05, 0.10\}$ $\times$ 3 random seeds).
The optimal configuration ($\eta=5.0$, $\gamma=0.10$) achieves 19.8\% lower regret compared to conservative settings.
Results exhibit low variance across random seeds (std=3.40), confirming algorithmic stability and reproducibility.

\paragraph{Convergence Guarantees.}
Log-log regression analysis reveals sublinear regret growth with exponent $\beta=0.669$ (R$^2$=0.9903), satisfying the PAC-learnability criterion ($\beta < 1$).
Cumulative regret grows as $O(T^{0.669})$, better than the worst-case linear bound $O(T)$ and approaching the theoretical $O(\sqrt{T})$ bound for contextual bandits, confirming efficient online learning.

\paragraph{Natural Cost Efficiency.}
The learned policy achieves \$2.43 per 1M tokens, a 75.7\% reduction compared to always using GPT-4-Turbo (\$10/1M).
Critically, the algorithm optimizes purely for quality ($\lambda_{\text{cost}} = 0$) without explicit cost penalties.
Cost efficiency emerges naturally from the model selection: GPT-4o (70.8\% usage) provides comparable quality to GPT-4-Turbo at 4$\times$ lower cost (\$2.50 vs \$10.00 per 1M tokens), and Mixtral (23.2\% usage) handles routine tasks effectively at \$0.27/1M.
This demonstrates that cost savings arise as a byproduct of quality-driven learning when the warmup expert's preference for expensive models is not justified by performance differences.

\paragraph{Practical Implications.}
These findings provide actionable guidance for deploying multi-model routing systems in production:

\textbf{When to Use Corralling:} Deploy the two-expert architecture when facing potential domain mismatch between training data and production workload. For example, if warmup priors come from general-purpose datasets (e.g., RouteLLM on LMSYS) but your application serves a specialized domain (e.g., legal, medical, customer support), Corralling provides automatic adaptation without manual retraining. The $O(T^{0.669})$ convergence rate means the system efficiently adapts within hundreds of samples.

\textbf{Cost Savings at Scale:} The 75.7\% cost reduction observed in this experiment illustrates the potential magnitude of routing savings. For example, at 1M tokens/day, switching from GPT-4-Turbo-only (\$10/1M) to the learned policy (\$2.43/1M) would save \$7.57/day. However, actual savings depend on the task mix, model pricing (which changes frequently), and the degree to which the deployment distribution matches the evaluation conditions. These projections should be treated as illustrative, not guaranteed.

\textbf{Handling New Models:} When new models become available (e.g., GPT-4o release), semantic transfer provides practical cold-start mitigation despite eventually being unlearned. Initialize the new model's priors from the most similar existing model (determined by cost tier, capability, or provider) with conservative gamma scaling ($\gamma=0.05$). This provides a reasonable starting point while allowing rapid online adaptation.

\textbf{Two Adaptation Regimes:}
\begin{itemize}[nosep]
    \item \textbf{Short-term (first $\sim$100--300 steps):} Semantic transfer provides immediate benefit via implicit regularization (breaking symmetry). This accelerates adoption compared to cold start, as demonstrated in Figure~\ref{fig:accelerated-adoption} with conservative learning ($\eta=0.1$).
    \item \textbf{Long-term ($>$500 steps with aggressive learning):} With sufficient data and higher learning rate ($\eta=5.0$), the system unlearns suboptimal priors and converges to the true optimal policy. Our results show the router discovered GPT-4o as the dominant choice (70.8\%) by adapting beyond the initial GPT-4-Turbo-biased semantic transfer.
\end{itemize}

\textbf{Design Implication:} The unlearning behavior is a \emph{feature, not a bug}. It ensures the system is not locked into potentially incorrect semantic priors while still benefiting from transfer during the critical cold-start phase. For production deployment, use conservative learning ($\eta \sim 0.1$--1.0) initially to exploit semantic priors, then increase learning rate ($\eta \sim 2.0$--5.0) after collecting sufficient data to enable adaptation.

\textbf{Hyperparameter Selection:} For deployment, start with $\eta=5.0$ and $\gamma=0.10$ as shown to be robust across seeds (std=3.40). If your application prioritizes safety over adaptation speed, use more conservative values ($\eta=1.0$, $\gamma=0.05$). The ablation study (Appendix~\ref{sec:appendix_corralling_ablation}) shows performance remains strong across $\eta \in [0.5, 5.0]$ when $\gamma \geq 0.05$, providing flexibility for domain-specific tuning.

\textbf{Sample Efficiency:} The system achieves strong performance with only N=1,121 training samples, making it practical for production deployment where labeled feedback is expensive. The sublinear regret growth ($\beta=0.669$) means performance improves steadily but with diminishing marginal gains---reaching 93.9\% quality after 1,121 samples suggests limited benefit from massive-scale data collection.

\textbf{Production Deployment:} The learned model distribution (GPT-4o: 70.8\%, Mixtral: 23.2\%, GPT-4-Turbo: 6.0\%) provides guidance for capacity planning and rate limit allocation. Most traffic concentrates on mid-tier models, enabling cost-efficient infrastructure design. Monitor expert weight evolution in production: if weights stabilize within 200-300 samples (as in Figure~\ref{fig:corralling_semantic}), the system has converged; if weights continue fluctuating, investigate potential distribution drift or data quality issues.


% ============================================================================
% Semantic Projection: Visualization on 1M Space
% ============================================================================

\subsection{Semantic Projection onto 1M Prompts}
\label{sec:semantic_projection}

\paragraph{Separation of Optimization and Visualization.}
A critical methodological point: we \emph{only} train on data where we have actual rewards (N=1,121 dev dataset samples).
We then \emph{project} the learned policy onto the full 1M semantic space to visualize coverage, but we do \emph{not} evaluate rewards on the 1M dataset (as we lack ground truth).

This ensures mathematical soundness:
\begin{itemize}[nosep]
    \item \textbf{Optimization Phase}: Learn on labeled data using importance-weighted loss (Eq.~\ref{eq:importance_weighted_loss})
    \item \textbf{Visualization Phase}: Project learned parameters onto 1M semantic space to show cluster structure
\end{itemize}

\paragraph{Cluster Analysis.}
We embed all 1M prompts using Sentence-BERT~\cite{reimers2019sentence} and project onto the first two principal components.
We identify two clusters based on projection analysis (N=50,000 sample):
\begin{itemize}[nosep]
    \item \textbf{Easy Cluster} (PC1 $< 0.3$): 47,085 prompts (94.2\%)
    \item \textbf{Hard Cluster} (PC1 $\geq 0.3$): 2,915 prompts (5.8\%)
\end{itemize}

The Easy cluster exhibits high semantic density (visible in KDE contours, Figure~\ref{fig:corralling_semantic} left) and contains prompts that can be served effectively by cheaper models.
This validates our hypothesis that the semantic structure enables cost-quality optimization.
The cluster distribution remains stable across different sample sizes, confirming the robustness of the semantic manifold structure.

\paragraph{Policy Projection.}
For each point in the 1M space, we determine which model the learned Corralling policy would select:
\begin{enumerate}[nosep]
    \item Sample expert $e \sim p$ according to learned weights
    \item Query expert $e$ for model selection: $a = \text{Expert}_e(x)$
    \item Record selection (no reward evaluation)
\end{enumerate}

This shows the \emph{coverage} of the learned policy across the semantic manifold without requiring ground truth rewards.

\paragraph{Key Insight.}
The Easy cluster (94.2\% of prompts) can be served with comparable quality by mid-tier models, which Corralling discovers automatically by unlearning the warmup bias.
This demonstrates that semantic structure enables generalization: by learning on 1,121 labeled samples, we can infer policy behavior across the full 1M distribution.
The cost savings (75.7\% reduction) are not from explicit cost optimization, but from correcting a quality-based bias---the warmup expert's false belief that expensive flagship models are necessary for high quality on routine tasks.


% ============================================================================
% References (add to main bibliography)
% ============================================================================

% @inproceedings{agarwal2017corralling,
%   title={Corralling a band of bandit algorithms},
%   author={Agarwal, Alekh and Luo, Haipeng and Neyshabur, Behnam and Schapire, Robert E},
%   booktitle={Conference on Learning Theory},
%   pages={12--38},
%   year={2017},
%   organization={PMLR}
% }

% @inproceedings{reimers2019sentence,
%   title={Sentence-BERT: Sentence embeddings using Siamese BERT-networks},
%   author={Reimers, Nils and Gurevych, Iryna},
%   booktitle={Proceedings of the 2019 Conference on Empirical Methods in Natural Language Processing},
%   year={2019}
% }

